\section{결과}
% 이 절의 수치는 모두 analysis/run_analysis.py가 만든 매크로다. 숫자를 손으로 고치지 않는다.

\subsection{참가자와 시행}
수집 기간에 서버에 도착한 완료 세션은 \Nrecruited{}개였다.
사전 등록한 기준으로 \Nexcluded{}개를 제외했다.
사유는 파일럿 \NexclPilot{}개, 수집 기간 밖 \NexclWindow{}개, 재참여 \NexclRepeat{}개, 미완료 \NexclIncomplete{}개, 유효 시행 부족 \NexclFewValid{}개였다.
남은 \Nanalyzed{}명의 목표 시행 \TrialsAnalyzed{}회 중 \TrialsExcluded{}회(\TrialsExcludedPct\%)가 유효 시행 정의를 만족하지 못해 빠졌다.
참가자별 유효 시행 수의 조건 간 차이는 평균 \Imbalance{}회였다.
분석한 참가자의 주 입력 방식은 터치 \NinputTouch{}명, 마우스 \NinputMouse{}명, 키보드 \NinputKey{}명, 혼합 \NinputMixed{}명이었다.
기록된 경고--목표 간격은 \LagMin--\LagMax\,ms, 테두리 점등 시간은 \CueOnMin--\CueOnMax\,ms였다.
조작 점검에서 테두리를 봤다는 응답은 \SawYes{}명, 못 봤다 \SawNo{}명, 모르겠다 \SawUnsure{}명이었다.
제외된 세션과 시행의 사유는 재현 패키지의 \texttt{analysis/output/exclusions.csv}와 \texttt{results.json}에 있다.

\subsection{기술통계}
표~\ref{tab:desc}는 참가자별 조건 중앙값을 다시 참가자들에 걸쳐 요약한 것이다.
그림~\ref{fig:paired}은 참가자 한 명을 선 하나로 이어, 두 조건의 값을 나란히 보여 준다.

\begin{table}[h]
\centering\small
\caption{참가자별 반응시간 중앙값의 요약 (단위 ms, $n=\Nanalyzed$)}\label{tab:desc}
\IfFileExists{generated/table_desc.tex}{\input{generated/table_desc}}{\Pend}
\end{table}

\begin{figure}[h]
\centering
\IfFileExists{figures/fig_paired.pdf}{\includegraphics[width=0.48\linewidth]{figures/fig_paired.pdf}}{\fbox{\Pend}}\hfill
\IfFileExists{figures/fig_diff.pdf}{\includegraphics[width=0.48\linewidth]{figures/fig_diff.pdf}}{\fbox{\Pend}}
\caption{왼쪽: 참가자별 조건 중앙값. 회색 선 하나가 참가자 한 명이고, 검은 선은 참가자 중앙값이다.
오른쪽: 참가자별 차이(신호 있음 $-$ 없음). 점 하나가 참가자 한 명이고, 세로선은 Hodges--Lehmann 추정치, 음영은 그 95\% 신뢰구간이다.
0보다 왼쪽이면 신호가 있을 때 빨랐다는 뜻이다.}\label{fig:paired}
\end{figure}

\subsection{확증 분석}
참가자별 차이 $d$의 Hodges--Lehmann 추정치는 \HL\,ms(95\% CI [\HLCIlo, \HLCIhi], 단측 95\% 상한 \HLupper)였다.
사전 등록한 단측 Wilcoxon 부호순위 검정 결과는 $W^{+}=\Wplus$, $p$ \Wp{}였다.
사전에 정한 판정 규칙에 따라 가설은 \textbf{\ResVerdict}.
참고로 $d$의 표본 중앙값은 \DiffMedian\,ms였다.
이 값은 위에 적은 조건별 참가자 중앙값(신호 있음 \MedCue\,ms, 없음 \MedNocue\,ms)의 차와 같지 않다.
중앙값의 차는 차의 중앙값과 일반적으로 다르기 때문이다.

\subsection{보조 분석과 민감도 분석}
평균 차이는 \DiffMean\,ms(표준편차 \DiffSD, 95\% CI [\DiffCIlo, \DiffCIhi], 단측 95\% 상한 \DiffUpper)였다.
단측 대응 $t$검정은 $t(\Tdf)=\Tstat$, $p$ \Tp{}였고, 효과크기는 $d_z=\Dz$였다.
차이가 0이 아닌 \Nnonzero{}명 중 \Kfaster{}명이 신호가 있을 때 더 빨랐고, \Kslower{}명은 더 느렸다(부호 검정 $p$ \SignP).
반응시간 하한을 150\,ms로 올리면($n=\SensFloorN$) HL 추정치는 \SensFloorHL\,ms, $p$ \SensFloorP{}였다.
중앙값 대신 10\% 절사평균을 쓰면 HL 추정치는 \SensTrimHL\,ms, $p$ \SensTrimP{}였다.
이 표본 크기에서 단측 대응 $t$검정이 검정력 .80으로 잡아낼 수 있는 최소 효과크기는 $d_z\approx\MDEdz$이다(주 검정인 Wilcoxon 기준이 아니다).

\subsection{전략 점검}
캐치 시행 오반응은 신호 있음 조건에서 \FAcue{}건(\FAcuePct\%), 신호 없음 조건에서 \FAnocue{}건(\FAnocuePct\%)이었다.
참가자별 오반응 수의 조건 차이에 대한 양측 부호순위 검정은 $p$ \FAdiffP{}였다.
조기 반응은 신호 있음 \AntCue{}건, 신호 없음 \AntNocue{}건이었다(조건별 목표 시행 \TrialsPerCond{}회 중).
100--150\,ms 반응의 비율은 신호 있음 \FastCuePct\%, 신호 없음 \FastNocuePct\%였다.

\subsection{탐색 분석 (사전 등록한 확증 분석이 아님)}
입력 방식별, 조작 점검 응답별 차이는 \texttt{results.json}의 \texttt{exploratory}에 있다.
집단별 인원이 적어 해석하지 않는다.

\TODO{데이터를 본 뒤 작성: 가설대로 나온 부분과 다르게 나온 부분을 나눠 한 문단으로 정리 (T10-C32, C35). 반대 방향 참가자 수(\Kslower)와 그 크기, 전략 점검이 가리키는 바를 숨기지 말고 적는다.}
